\chapter{Future Work}

Methdological extension that could increase the planning:
\begin{itemize}
    \item Origin-destination traffic flow data on sub-annual level
    \item Increasing robustness in modeling algorithms considering uncertainties
    \item Cross-modal consideration while keeping the geographic scope and increasing geographic resolution 
\end{itemize}

Concrete research questions to address with these methodologies under the \textit{geographic allocation of charging infrastructure capacities}:
\begin{itemize}
    \item Going from covering demand, to support BEV adoption, extending the horizon beyond boarders and understanding 2nd-hand adoption dynamics and how these could intertwine with the investments there 
    \item understanding the willingness to adapt logistics to the cost-optimal BEV operation of fleet operators (and passenger cars)? (i.e., going beyond demand-response, but with behavioral flexibility) particularly interesting for fleet operators;
    \item Increasing the intersection between energy supply to multiple charging/fueling infrastructure; going beyond consumer-overarching and road vehicle/segment-type overarching but to also consider mode-overarching in long-haul and international freight transport to consider for these synergies in the energy supply, while still keeping the geographic scope;
\end{itemize}






